\subsection{Etap Code}
Faza wytwarzania (\textit{Code}) obejmuje implementację wymagań przy zachowaniu standardów jakości kodu i~spójności historii zmian. Fundamentem tego etapu jest rozproszony system kontroli wersji Git, a~konwencje nazewnictwa zatwierdzeń są wymuszane na poziomie procesu CI.

Proces kodowania realizowany jest przy użyciu nowoczesnych środowisk programistycznych (\textit{IntelliJ IDEA}, \textit{VS Code}), a~sama implementacja opiera się na technologiach realizujących paradygmat mikroserwisów.

Wszystkie komponenty systemu są projektowane zgodnie z~zasadami \textit{Twelve-Factor App}. Kluczowe aspekty to:
\begin{itemize}
    \item \textbf{Konfiguracja:} Cała konfiguracja specyficzna dla środowiska jest wstrzykiwana przez zmienne środowiskowe, a~nie zapisana na stałe w~kodzie.
    \item \textbf{Procesy:} Aplikacje są bezstanowe (ang. \textit{stateless}), co umożliwia ich łatwe skalowanie w~poziomie wewnątrz klastra Kubernetes.
    \item \textbf{Logi:} Są traktowane jako strumień zdarzeń wysyłany na standardowe wyjście, co umożliwia ich centralizację przez narzędzia obserwowalności.
\end{itemize}


\subsubsection{Backend (Java + Quarkus + Native Image)}
Mikroserwisy backendowe zbudowano z~użyciem frameworka \textit{Quarkus} i~skompilowano do natywnych plików wykonywalnych za pomocą \textit{GraalVM Native Image} (kompilacja \textit{Ahead-of-Time}, AOT). W~trybie natywnym Quarkus analizuje ścieżki wykonania kodu w~czasie kompilacji i~wyklucza nieużywane klasy, co skutkuje:
\begin{itemize}
    \item \textbf{Błyskawicznym startem} --- aplikacja jest gotowa do obsługi żądań w~czasie poniżej 100 ms, co skraca czas wznowienia po awarii i~przyspiesza skalowanie poziome w~Kubernetes.
    \item \textbf{Niskim zużyciem pamięci} --- brak środowiska uruchomieniowego JVM zmniejsza ślad pamięciowy kontenera, co bezpośrednio przekłada się na efektywność kosztową w~środowiskach chmurowych.
    \item \textbf{Minimalnym obrazem kontenera} --- obraz oparty na \texttt{ubi-minimal} zawiera wyłącznie system operacyjny i~skompilowany plik binarny, bez JDK ani bibliotek środowiska uruchomieniowego.
\end{itemize}
Wstrzykiwanie zależności realizowane jest przez mechanizm CDI Quarkusa (\textit{ArC}), który w~trybie AOT przetwarza wszystkie adnotacje w~czasie kompilacji, a~nie w~czasie uruchomienia.

\begin{techcard}{Quarkus}

\textbf{Opis:} Quarkus to nowoczesny stos technologiczny Java, określany jako \textit{Supersonic Subatomic Java}. Zaprojektowany z~myślą o~środowiskach \textit{Cloud Native} i~kontenerach, charakteryzuje się bardzo niskim zużyciem pamięci i~błyskawicznym czasem uruchamiania, co jest kluczowe w~architekturze mikrousługowej.

\textbf{Dobre nawyki:}
\begin{itemize}
  \item \textbf{Testowanie natywne (\texttt{@QuarkusIntegrationTest}):} Testy w~trybie natywnym weryfikują, czy optymalizacje AOT nie usunęły elementów wymaganych przez refleksję.
  \item \textbf{Dev Services:} Automatyczne uruchamianie baz danych i~brokerów w~kontenerach podczas prac programistycznych zapewnia spójność środowisk u~wszystkich programistów.
  \item \textbf{Umiar w~refleksji:} Należy preferować rozwiązania obsługujące indeksowanie Jandex i~unikać bibliotek intensywnie korzystających z~dynamicznego ładowania klas.
  \item \textbf{Continuous Testing:} Tryb ciągłego testowania w~konsoli zapewnia natychmiastową informację zwrotną podczas pisania kodu.
\end{itemize}

\textbf{Rola w~projekcie:} Quarkus służy do budowy wydajnych mikroserwisów o~wysokiej gęstości upakowania w~klastrze Kubernetes. Aplikacje zużywają minimalną ilość zasobów, co przekłada się na niższe koszty infrastruktury. Konfigurację aplikacji prezentuje listing \ref{lst:quarkus-app-properties}, a~konfigurację natywnego buildu w~Gradle --- listing \ref{lst:quarkus-native-build} (etap Build).

\end{techcard}

\begin{lstlisting}[caption={Fragment konfiguracji aplikacji (apps/users/src/main/resources/application.properties)}, label={lst:quarkus-app-properties}]
quarkus.application.name=users
quarkus.http.port=8080

quarkus.datasource.db-kind=postgresql
quarkus.hibernate-orm.schema-management.strategy=update

mp.messaging.outgoing.user-events.connector=smallrye-pulsar
mp.messaging.outgoing.user-events.topic=persistent://public/default/user-events
\end{lstlisting}

Zwraca uwagę brak jawnych danych uwierzytelniających w~pliku konfiguracyjnym: w~środowisku deweloperskim ich nieobecność aktywuje mechanizm Dev Services, który automatycznie uruchamia bazę PostgreSQL w~kontenerze, natomiast w~klastrze Kubernetes dane te wstrzykiwane są przez zmienne środowiskowe, zgodnie z~zasadami \textit{Twelve-Factor App}.


\subsubsection{Frontend (Vue 3 + Vite)}
Zgodnie z~decyzją \textit{ADR-0014}, warstwę prezentacji zbudowano z~użyciem frameworka \textit{Vue 3} i~narzędzia budującego \textit{Vite}. Kod pisany jest w~\textit{TypeScript}, a~poprawność typów weryfikuje \textit{vue-tsc} przed każdym buildem produkcyjnym. Komponenty interfejsu komunikują się z~mikroserwisami backendowymi przez REST API.

\begin{techcard}{Vue~3}

\textbf{Opis:} Vue~3 to progresywny framework JavaScript do budowania interfejsów użytkownika. Wprowadza \textit{Composition API} --- umożliwiające organizowanie logiki komponentu według funkcjonalności. Projekt kompilowany jest narzędziem Vite, które wykorzystuje natywne moduły ES podczas prac programistycznych i~Rollup podczas produkcyjnego budowania.

\textbf{Dobre nawyki:}
\begin{itemize}
  \item \textbf{TypeScript przez \texttt{vue-tsc}:} Weryfikacja typów w~komponentach \texttt{.vue} przed buildem produkcyjnym bez generowania plików wyjściowych.
  \item \textbf{Proxy Vite:} Transparentne przekierowanie zapytań do \texttt{/api} na serwer backendowy eliminuje potrzebę konfigurowania CORS w~środowisku deweloperskim.
  \item \textbf{Multi-stage w~Dockerze:} Etap \texttt{npm ci \&\& vite build} produkuje statyczne pliki; obraz docelowy zawiera wyłącznie \texttt{nginx} bez środowiska Node.js.
\end{itemize}

\textbf{Rola w~projekcie:} Vue~3 pełni rolę cienkiej warstwy prezentacji (ang. \textit{thin client}) nad REST API mikroserwisów. Warstwa frontendowa wdrażana jest jako osobny mikroserwis w~Kubernetes z~własnym procesem CI/CD, analogicznie do serwisów backendowych.

\end{techcard}

\begin{lstlisting}[caption={Fragment konfiguracji Vite z~proxy dla środowiska deweloperskiego (apps/frontend/vite.config.ts)}, label={lst:vite-config}]
import { defineConfig } from 'vitest/config';
import vue from '@vitejs/plugin-vue';

export default defineConfig({
  plugins: [vue()],
  server: {
    proxy: {
      '/api/users': {
        target: 'http://localhost:8080',
        rewrite: (path) => path.replace(/^\/api\/users/, '/users'),
      },
      '/api/tasks': {
        target: 'http://localhost:8080',
        rewrite: (path) => path.replace(/^\/api\/tasks/, ''),
      },
    },
  },
});
\end{lstlisting}


\subsubsection{System kontroli wersji: GitHub}
Jako system kontroli wersji i~platformę współpracy wybrano \textbf{GitHub}, oparty na rozproszonym systemie Git. Jest on jedynym źródłem prawdy (ang. \textit{single source of truth}) dla całego monorepozytorium.

\textbf{Standard Conventional Commits:}
Zgodnie z~decyzją projektową \textit{ADR-0004}, wprowadzono standard \textbf{Conventional Commits}. Polega on na nadawaniu opisom zatwierdzeń ustrukturyzowanej formy (np. \texttt{feat:}, \texttt{fix:}, \texttt{chore:}) \cite{conventionalCommits}.
\begin{itemize}
    \item \textbf{Uzasadnienie:} Standaryzacja pozwala na automatyczne generowanie list zmian (ang. \textit{changelogs}) oraz automatyzację wersjonowania semantycznego (ang. \textit{Semantic Versioning}).
    \item \textbf{Rola w~CI/CD:} Czytelna historia zmian umożliwia narzędziom automatyzującym precyzyjne określenie, które komponenty systemu wymagają przebudowania i~wdrożenia po każdej zmianie w~monorepozytorium.
\end{itemize}

Poniższy przykład ilustruje kompletny opis zatwierdzenia zgodny ze standardem. Nagłówek zawiera typ, zakres oraz zwięzły opis zmiany, a~stopka (ang. \textit{footer}) powiązuje commit z~konkretnym zadaniem w~systemie Jira:

\begin{lstlisting}[caption={Przykład opisu zatwierdzenia zgodnego ze standardem Conventional Commits}, label={lst:conventional-commit}]
feat(helm): create umbrella chart for all 7 services

Aggregates all microservice charts under a single parent
chart, enabling one-command deployment of the full stack.

Resolves: INZ-9
\end{lstlisting}

